\scanpage{150}
\unsure{partie|partir} de $***$. Mais, si \deleted{$u\to u_{\unsure{0|\infty}}$}\inserted{$1$ n'est pas valeur propre,} on a aussi, \unsure{\deleted{pour \unsure{tout|tt}}|pour \unsure{tout|tt}} \inserted{\unsure{pour la dimension infinie|particulièrement en dimension infinie}} $\unsure{\rho|p}_i\deleted{\unsure{\geq}}=\frac{\det(1+\unsure{\lambda|u})}{1+\lambda_i}\deleted{\unsure{\geq|>}}$ d'où

\[
\boxed{\|R(u)\|_\infty\geq|\det(1+u)|}\quad\text{(op normal)}
\]

(car il existe des $\lambda_i$ arbitrairement petits, donc des $\unsure{\rho|p}_i$ arbitrairement \deleted{\unsure{distants|inférieurs|inférieur} \unsure{par|que}}\inserted{voisins} de $\det(1\unsure{+|-}u)$ !

On \deleted{\illegible{1}} \unsure{prouve} donc que l'inégalité $***$ a) ne peut s'améliorer, car elle devient égalité pour $u\geq0$ !
